% !TEX root = ../thesis.tex
\chapter*{DANH MỤC KÝ HIỆU VÀ CHỮ VIẾT TẮT}
\addcontentsline{toc}{chapter}{DANH MỤC KÝ HIỆU VÀ CHỮ VIẾT TẮT}

\begin{longtable}{|>{\raggedright\arraybackslash}p{3.6cm}|>{\raggedright\arraybackslash}p{10.4cm}|}
  \hline
  \textbf{Chữ viết tắt} & \textbf{Nghĩa đầy đủ / Giải thích} \\
  \hline
  API   & Application Programming Interface (Giao diện lập trình ứng dụng) \\ \hline
  BERT  & Bidirectional Encoder Representations from Transformers \\ \hline
  BERTScore & Chỉ số đánh giá tương đồng ngữ nghĩa dựa trên vector ngữ cảnh BERT \\ \hline
  GUID  & Globally Unique Identifier (Mã định danh duy nhất toàn cầu) \\ \hline
  LLM   & Large Language Model (Mô hình ngôn ngữ lớn) \\ \hline
  MMR   & Maximal Marginal Relevance --- Thuật toán chọn lọc nội dung nổi bật và giảm dư thừa \\ \hline
  NLP   & Natural Language Processing (Xử lý ngôn ngữ tự nhiên) \\ \hline
  ROUGE & Recall-Oriented Understudy for Gisting Evaluation --- Bộ độ đo trùng lặp n-gram \\ \hline
  Seq2Seq & Sequence-to-Sequence (Mô hình biến đổi chuỗi sang chuỗi) \\ \hline
  SFT   & Supervised Fine-Tuning (Huấn luyện tinh chỉnh có giám sát) \\ \hline
  ViT5  & Vietnamese Text-to-Text Transfer Transformer (Mô hình T5 cho tiếng Việt) \\ \hline
  \multicolumn{2}{l}{} \\
  \hline
  \textbf{Thuật ngữ} & \textbf{Định nghĩa} \\ \hline
  Abstractive Summarization & Tóm tắt trừu tượng hóa: Sinh câu mới diễn giải lại nội dung văn bản gốc \\ \hline
  Extractive Summarization  & Tóm tắt trích xuất: Trích chọn nguyên văn các câu quan trọng từ bài gốc \\ \hline
  Selective Long-Context   & Kỹ thuật chọn lọc các đoạn ngữ cảnh quan trọng rải rác trong toàn văn bài báo dài \\ \hline
  Sentence-aware Chunking  & Phương pháp chia đoạn văn bản dựa theo ranh giới câu chuẩn xác \\ \hline
  Overlapping Chunks        & Kỹ thuật chia đoạn có khoảng đè câu nhằm bảo toàn tính liên kết ngữ cảnh \\ \hline
  Keyword-conditioned Gen.  & Cơ chế điều kiện hóa quá trình sinh văn bản của mô hình bằng danh sách từ khóa \\ \hline
  Factuality Reranker       & Bộ xếp hạng lại ứng viên tóm tắt nhằm tối ưu độ chính xác thực thể và số liệu \\ \hline
  Hallucination             & Hiện tượng ảo giác: Mô hình sinh ra các thông tin sai lệch không có trong nguồn \\ \hline
\end{longtable}
